\section{Introduction}

This report is the result of a university course project: a hands-on study of the
socket API that Linux offers to a program that talks to many peers over TCP. The
goal was to learn these interfaces by using them, and then to measure what each
choice changes when the program is put under load. The parts of the API exercised
along the way include:
\begin{itemize}
  \item blocking and non-blocking sockets, \code{connect}, \code{accept},
        \code{send} and \code{recv}, and socket options such as \code{SO\_SNDTIMEO};
  \item readiness multiplexing with \code{poll(2)} and \code{epoll(7)}~\cite{poll},
        and waking a blocked poller with \code{eventfd(2)};
  \item ending a connection with \code{shutdown} and \code{close}, and what happens to
        a descriptor's number afterwards;
  \item process accounting with \code{getrusage(2)} and \code{/proc}, and network
        emulation with network namespaces, veth pairs and netem~\cite{netem}.
\end{itemize}

The vehicle for the study is a small peer-to-peer chat node written in C++23. The
node is built so that everything above the transport stays the same: conversations,
identities, the handshake protocol, the write policy and the user interface. Only
the transport changes, and it is chosen at start-up with \code{-{}-io=}. Nine
transports were written, each using the socket API in a different way. Because the
rest of the program is shared, a difference between two runs can be traced to how
the sockets are read and written.

\subsection{The three axes}

The nine backends vary three properties:
\begin{description}
  \item[Concurrency model] a thread per connection, a thread per peer, one shared
        reactor thread, or one reactor for reads and another for writes.
  \item[Readiness API] blocking calls, \code{poll(2)}, or \code{epoll(7)}~\cite{poll}.
  \item[Socket topology] \emph{simplex}, with two TCP connections per peer (one per
        direction), or \emph{duplex}, with one connection carrying both directions.
\end{description}

Table~\ref{tab:pairs} lists the pairs of modes that differ in a single factor. The
results in Section~\ref{sec:results} are read through these pairs.

\begin{table}[h]
\centering\small
\begin{tabular}{@{}ll@{}}
\toprule
Pair & Factor isolated \\
\midrule
\mode{poll} vs \mode{epoll} (and their duplex and split versions) & readiness API \\
\mode{threads} vs \mode{pool-sender} & blocking vs pooled non-blocking writes \\
\mode{poll}/\mode{epoll} vs \mode{*-duplex} & two connections per peer vs one \\
\mode{*-duplex} vs \mode{*-split} & one I/O thread vs separate read and write threads \\
\mode{threads} vs \mode{peer-threads} & thread per connection: two sockets vs one \\
\mode{peer-threads} vs \mode{poll-duplex} & thread per peer vs one shared poller \\
\bottomrule
\end{tabular}
\caption{Mode pairs that each isolate one factor.}
\label{tab:pairs}
\end{table}

\subsection{Research questions}

\begin{enumerate}[label=\textbf{Q\arabic*}]
  \item How do thread-per-connection modes compare with one I/O thread as peers, idle
        connections and slow consumers grow, and where does each one break?
  \item Does \code{poll}'s per-wake-up cost, which grows with the number of
        connections, become visible?
  \item Does a slow peer stall everyone when writes are blocking, and do pooled writes
        with drop-on-full isolate it?
  \item What does one connection per peer save, compared with two, in fds, handshake
        round trips and latency?
  \item Does flushing writes on a separate thread help?
  \item How does a thread per peer on one socket compare with \mode{threads} and with
        a shared poller?
\end{enumerate}

\subsection{Structure of the report}

\begin{itemize}
  \item Section~\ref{sec:arch} describes the part of the node that is the same in
        every mode, and the problems met while building it.
  \item Section~\ref{sec:wrappers} describes the C++ classes that wrap the socket,
        \code{poll} and \code{epoll} calls, and how they use the C functions.
  \item Section~\ref{sec:backends} describes the nine transports built from them.
  \item Section~\ref{sec:method} describes how the node was measured: its counters,
        the load generator, the test harness and the scenarios.
  \item Section~\ref{sec:results} presents the measurements scenario by scenario.
  \item Section~\ref{sec:compare} compares blocking calls, \code{poll} and
        \code{epoll} across all scenarios.
  \item Section~\ref{sec:conclusions} draws conclusions and ends with an untested
        design idea.
\end{itemize}
